% ==============================================================================
% SLIDE 8: INFRASTRUCTURE ADAPTATION: CABLE CAR & SOCIO-ECONOMIC HARMONY
% ==============================================================================
\begin{frame}[label=cablecar]{Environmental Adaptation: The Santorini Cable Car}
  \begin{columns}[T, onlytextwidth]
    % Left Column: Photo
    \begin{column}{0.48\textwidth}
      \slidephoto{Santorini-cable-car-station.jpg}{Fira Cable Car Upper Station: Aerial tramway conquering the 220\,m vertical caldera cliff}
      % Alternative choice: \slidephoto{cable-car.jpg}{Santorini Cable Car: Doppelmayr tramway traversing the volcanic cliff face}
    \end{column}
    
    % Right Column: Bullet Points
    \begin{column}{0.48\textwidth}
      \textbf{\color{aegeanblue}Engineering \& Socio-Economic Harmony:}
      \vspace{0.1em}
      \begin{itemize}
        \setlength{\itemsep}{0.12em}
        \item \textbf{The Caldera Bottleneck:} Ascending 220\,m of sheer volcanic cliff from the Old Port to Fira required tackling 587 steep zigzagging cobblestone steps.
        \item \textbf{Aerial Adaptation (Doppelmayr 1982):} An aerial cableway spans the 220\,m drop in 3 minutes (1,200 pass/hr), avoiding unstable cliff excavation.
        \item \textbf{The Nomikos Benefaction:} Shipowner Evangelos Nomikos funded 100\% via the Loulas Nomikos Foundation, donating it to the 14 communities.
        \item \textbf{The Muleteer Royalty Accord:} To protect local livelihoods, a statutory share of every ticket is distributed to traditional mule drivers.
      \end{itemize}
    \end{column}
  \end{columns}
\end{frame}
